\section{Discussion and Conclusion}\label{sec:discussion}

When aggregate benchmarks show equivalent performance, the choice among hybrid architectures is often driven by efficiency considerations. Per-example multiplicity reveals a complementary dimension: each backbone handles a different subset of inputs correctly, determined by the recurrent mechanism rather than the attention configuration (\S\ref{ssec:decoupling}). For deployment contexts where per-example reliability matters, aggregate equivalence provides incomplete assurance. The complementarity structure (Table~\ref{tab:complementarity}) shows that oracle-level combination yields $+7$--$14$\,pp gains, but majority vote captures only ${\sim}51$\% because majority-exploitable and majority-hostile disagreements occur in near-equal proportion, motivating architecture-aware routing beyond simple ensembling. Our within-family design isolates backbone variation as the source of this multiplicity; \citet{yang2026benchmark} document analogous per-item disagreement (16--66\%) among frontier LLMs under conflated training regimes, suggesting the phenomenon is partly architectural in origin and resists homogenization even at minimal attention ratios.
